\subsection{MAPPINGS COMPATIBLE WITH EQUIVALENCE RELATIONS}

Let R be an equivalence relation on a set E, and let f be a function with domain E. Then f is said to be compatible with the relation R if the relation \(y = f(x)\) is compatible (with respect to x) with the relation \(R\{x, x'\}\) .

It comes to the same thing to say that the restriction of \(f\) to each equivalence class is a constant mapping, in which case we say that \(f\) is constant on each equivalence class with respect to \(\mathbf{R}\) . If \(g\) is the canonical mapping of \(\mathbf{E}\) onto \(\mathbf{E} / \mathbf{R}\) , an equivalent statement is that the relation \(g(x) = g(x')\) implies \(f(x) = f(x')\) ; hence (§3, Proposition 9) we have the following criterion:

C57. Let R be an equivalence relation on a set E, and let g be the canonical mapping of E onto E/R. Then a mapping f of E into F is compatible with R if and only if f can be put in the form h o g, where h is a mapping of E/R into F. The mapping h is uniquely determined by f; if s is any section of g, we have \(h = f \circ s\) .

The mapping h is said to be induced by f on passing to the quotient with respect to R.

Let \(f\) be a mapping of a set \(\mathbf{E}\) into a set \(\mathbf{F}\) , and let \(\mathbf{A} = f\langle \mathbf{E}\rangle \subset \mathbf{F}\) . Let \(\mathbf{R}\) be the equivalence relation associated with \(f\) (no. 2); it is clear that \(f\) is compatible with \(\mathbf{R}\) . Moreover, the mapping \(h\) induced by \(f\) on passing to the quotient is an injection of \(\mathbf{E} / \mathbf{R}\) into \(\mathbf{F}\) : for if \(t, t'\) are equivalence classes with respect to \(\mathbf{R}\) such that \(h(t) = h(t')\) , we have \(f(x) = f(x')\) for all \(x \in t\) and \(x' \in t'\) , which implies \(t = t'\) by the definition of \(\mathbf{R}\) . Let \(k\) be the mapping of \(\mathbf{E} / \mathbf{R}\) onto \(\mathbf{A}\) which has the same graph as \(h\) ; then \(k\) is bijective. If \(j\) is the canonical injection of \(\mathbf{A}\) into \(\mathbf{F}\) and if \(g\) is the canonical mapping of \(\mathbf{E}\) onto \(\mathbf{E} / \mathbf{R}\) , then we may write

\[
f = j \circ k \circ g.
\]

This relation is called the canonical decomposition of \(f\) .

Let \(f\) be a mapping of a set \(E\) into a set \(F\) , let \(R\) be an equivalence relation on \(E\) , and let \(S\) be an equivalence relation on \(F\) . Let \(u\) be the canonical mapping of \(E\) onto \(E / R\) , and let \(v\) be the canonical mapping of F onto F/S. The mapping f is said to be compatible with the equivalence relations R and S if \(v \circ f\) is compatible with R; this means that the relation \(x \equiv x' \pmod{R}\) implies \(f(x) \equiv f(x') \pmod{S}\) . The mapping h of E/R into F/S induced by \(v \circ f\) on passing to the quotient with respect to R is then called the mapping induced by f on passing to the quotients with respect to R and S; it is characterized by the relation \(v \circ f = h \circ u\) .

